% Diagramme de cas d'utilisation
\begin{figure}[H]
\centering
\begin{tikzpicture}[node distance=8mm and 18mm, font=\small]
  % Acteur principal a gauche
  \node (recruteur) at (0,-3) {\ucactorperson{Recruteur}};

  % Cas d'utilisation au centre dans la boite systeme
  \node[ucase] (uc1) at (5,0)    {Lancer un sourcing};
  \node[ucase] (uc2) at (5,-1.4) {Suivre un run\\en direct};
  \node[ucase] (uc3) at (5,-2.8) {Valider / éditer\\la sélection (HITL)};
  \node[ucase] (uc4) at (5,-4.2) {Consulter un\\rapport};
  \node[ucase] (uc5) at (5,-5.6) {Explorer la\\mémoire RAG};
  \node[ucase] (uc6) at (5,-7.0) {Consulter les\\métriques};

  % Boite systeme englobant les cas
  \begin{scope}[on background layer]
    \node[ucsystem, fit=(uc1)(uc2)(uc3)(uc4)(uc5)(uc6),
          label={[font=\small\bfseries, anchor=south]north:Système SMA de recrutement}] {};
  \end{scope}

  % Acteurs secondaires a droite
  \node (ollama) at (10.8,-1.0) {\ucactorperson{LLM Ollama}};
  \node (ddg)    at (10.8,-3.0) {\ucactorperson{DuckDuckGo}};
  \node (github) at (10.8,-5.0) {\ucactorperson{GitHub API}};
  \node (so)     at (10.8,-7.0) {\ucactorperson{Stack Exchange}};

  % Liens recruteur -> cas
  \draw (recruteur) -- (uc1);
  \draw (recruteur) -- (uc2);
  \draw (recruteur) -- (uc3);
  \draw (recruteur) -- (uc4);
  \draw (recruteur) -- (uc5);
  \draw (recruteur) -- (uc6);

  % Liens cas -> acteurs externes
  \draw (uc1) -- (ollama);
  \draw (uc1) -- (ddg);
  \draw (uc1) -- (github);
  \draw (uc1) -- (so);
  \draw (uc3) -- (ollama);
  \draw (uc5) -- (ollama);

  % Relations include
  \draw[-{Stealth[length=2.5mm]}, dashed] (uc1) -- node[font=\scriptsize, sloped, above]{\textit{<<include>>}} (uc2);
  \draw[-{Stealth[length=2.5mm]}, dashed] (uc2) -- node[font=\scriptsize, sloped, above]{\textit{<<extend>>}} (uc3);
\end{tikzpicture}
\caption[Diagramme de cas d'utilisation]{Diagramme de cas d'utilisation. Le recruteur humain interagit avec six cas d'usage exposés par l'interface graphique. Le système mobilise quatre acteurs externes : le LLM Ollama (local) et trois sources publiques de \emph{sourcing}. La relation \emph{include} indique que tout lancement de \emph{run} ouvre la page de suivi ; \emph{extend} signale que la validation HITL n'apparaît que sur les \emph{runs} où elle est activée.}
\label{fig:cas_utilisation}
\end{figure}
